\subsection{APPLICATION TO INVERSE LIMITS}

PROPOSITION 5. Let \((\mathbf{G}_{\alpha}, f_{\alpha \beta})\) be an inverse system of Hausdorff topological groups such that the \(f_{\alpha \beta}\) are surjective strict morphisms with compact kernels. Then for each \(\alpha \in I\) the canonical mapping \(f_{\alpha}\) of \(\mathbf{G} = \varprojlim \mathbf{G}_{\alpha}\) into \(\mathbf{G}_{\alpha}\) is a surjective strict morphism whose kernel is compact.

The facts that \(f_{\alpha}\) is surjective and that its kernel is compact are consequences of Chapter I, § 9, no. 6, Corollary 1 to Proposition 8. It remains to show that \(f_{\alpha}\) is a strict morphism. Let \(e\) (resp. \(e_{\alpha}\) ) denote the identity element of G (resp. \(G_{\alpha}\) ). Each neighbourhood V of \(e\) in G contains a set of the form \(\overline{f}_{\beta}^{-1}(V_{\beta})\) , where \(V_{\beta}\) is a neighbourhood of \(e_{\beta}\) in \(G_{\beta}\) , and we may suppose that \(\beta \geqslant \alpha\) ; since \(f_{\alpha \beta}\) is a surjective strict morphism, \(f_{\alpha \beta}(V_{\beta})\) is a neighbourhood of \(e_{\alpha}\) in \(G_{\alpha}\) , and since \(f_{\beta}\) is surjective, we have \(V_{\beta} \subset f_{\beta}(V)\) , whence

\[
f _ {\alpha} (\mathbf {V}) = f _ {\alpha \beta} (f _ {\beta} (\mathbf {V})) \supset f _ {\alpha \beta} (\mathbf {V} _ {\beta});
\]

this shows that \(f_{\alpha}(\mathbf{V})\) is a neighbourhood of \(e_{\alpha}\) in \(\mathbf{G}_{\alpha}\) .

If \(\mathrm{H}_{\alpha}=\overline{f}_{\alpha}^{-1}(e_{\alpha})\) , then each \(H_{\alpha}\) is a compact normal subgroup of G; the \(H_{\alpha}\) clearly satisfy the condition (AP) of no. 3; and \(G_{\alpha}\) can be identified with \(G/H_{\alpha}\) . In particular Propositions 3 and 4 of no. 3 apply to G and the \(H_{\alpha}\) .

COROLLARY 1. Let \((\mathbf{G}_{\alpha}, f_{\alpha \beta})\) be an inverse system of topological groups which satisfies the hypotheses of Proposition 5; let \((\mathbf{G}_{\alpha}', f_{\alpha \beta}')\) be an inverse system of topological groups, and for each \(\alpha\) let \(u_{\alpha}: \mathbf{G}_{\alpha} \to \mathbf{G}_{\alpha}'\) be a surjective strict morphism with compact kernel, such that the \(u_{\alpha}\) form an inverse system of mappings. Then \(u = \varprojlim u_{\alpha}\) is a strict morphism of \(\mathbf{G} = \varprojlim \mathbf{G}_{\alpha}\) onto \(\mathbf{G}' = \varprojlim \mathbf{G}_{\alpha}'\) , and its kernel is compact.

Let \(N_{\alpha}\) be the kernel of \(u_{\alpha}\) . \(L_{\alpha} = \overline{f}_{\alpha}^{-1}(N_{\alpha})\) is then the kernel of the surjective strict morphism \(v_{\alpha} = u_{\alpha} \circ f_{\alpha}: G \to G'_{\alpha}\) ; since \(L_{\alpha}/H_{\alpha}\) is isomorphic to \(N_{\alpha}\) (§ 2, no. 7, Proposition 20), \(L_{\alpha}\) is a compact normal subgroup of G (§ 4, no. 1, Corollary 2 to Proposition 2). The kernel L of u is the intersection of the \(L_{\alpha}\) . Let \(\varphi\) denote the canonical mapping \(G \to G/L\) ; then we can write \(v_{\alpha} = w_{\alpha} \circ \varphi\) , where \(w_{\alpha}\) is a strict morphism of G/L onto \(G'_{\alpha}\) , whose kernel is \(L_{\alpha}/L\) . Since the intersection of the \(L_{\alpha}/L\) is the identity element of G/L, and since the \(L_{\alpha}/L\) form a filter base and are compact, this filter base converges to the identity element of G/L (Chapter I, § 9, no. 1, Corollary to Theorem 1). Proposition 2 of no. 3 then shows that \(w = \varprojlim w_{\alpha}\) is an isomorphism of G/L onto \(G'\) ; it follows that \(w \circ \varphi\) is a strict morphism of G onto \(G'\) , with kernel L. But it is clear that \(u = w \circ \varphi\) , and thus the corollary is proved.

COROLLARY 2. Let \((\mathbf{G}_{\alpha}, f_{\alpha \beta})\) be an inverse system of topological groups which satisfies the conditions of Proposition 5, and let \(\mathbf{G}'\) be a topological group in which there is a neighbourhood \(\mathbf{V}'\) of the identity element \(e'\) which contains no subgroup of \(\mathbf{G}'\) other than \(\{e'\}\) . Then if \(v: \mathbf{G} \to \mathbf{G}'\) is any continuous homomorphism, there is an index \(\alpha \in \mathbf{I}\) and a continuous homomorphism \(v_{\alpha}: \mathbf{G}_{\alpha} \to \mathbf{G}'\) such that \(v = v_{\alpha} \circ f_{\alpha}\) .

For since \(\overline{v}^{-1}(\mathbf{V}^{\prime})\) is a neighbourhood of \(e\) in \(\mathbf{G}\) , there is an index \(\alpha\) and a neighbourhood \(\mathbf{V}_{\alpha}\) of \(e_{\alpha}\) in \(\mathbf{G}_{\alpha}\) such that \(\overline{f}_{\alpha}^{-1}(\mathbf{V}_{\alpha}) \subset \overline{v}^{-1}(\mathbf{V}^{\prime})\) . Hence \(v(\mathbf{H}_{\alpha}) \subset \mathbf{V}^{\prime}\) , and since \(v(\mathbf{H}_{\alpha})\) is a subgroup of \(\mathbf{G}^{\prime}\) , it follows that \(v(\mathbf{H}_{\alpha}) = \{e^{\prime}\}\) . Since \(f_{\alpha}\) may be identified with the canonical mapping \(\mathbf{G} \to \mathbf{G}/\mathbf{H}_{\alpha}\) , the corollary is a consequence of the canonical factorization of a continuous homomorphism (§ 2, no. 8).

